\section{Introduction}
\label{sec:introduction}

Automated recognition of crackles and wheezes can reduce the manual effort
involved in respiratory sound assessment\cite{hsu2021hflung}. However, training
respiratory sound classifiers requires expert annotations that are costly to obtain at scale\cite{zhang2024opera}. Joint training across
existing public datasets\cite{rocha2019icbhi,zhang2022sprsound} therefore offers a way to expand supervision without collecting new expert annotations.
Those datasets' annotations refer to shared acoustic attributes but differ in granularity and prediction targets, so their labels are not directly interchangeable.

Existing methods include mapping diverse annotations from multiple datasets to a unified prediction target to align task definitions. For instance, fusion experiments in SPRSound merged source categories \cite{zhang2022sprsound}, whereas LungMix mapped the ICBHI, SPRSound, and HF datasets to the same four-class label space used by ICBHI to facilitate single-source transfer \cite{ge2025lungmix}.
Multi-breath and PC-MCL express sound composition through multilabel supervision, but their training and evaluation remain within
ICBHI\cite{chua2024multibreath,jeong2026pcmcl}. DCASE supports joint training with class mappings and missing-label masks while retaining dataset-specific evaluation but both datasets share a temporal sound-event detection output
interface\cite{cornell2024dcase}. 

In contrast, our respiratory sound analysis task simultaneously performs the four-class respiratory cycle detection from the ICBHI dataset and the binary classification task from the SPRsound dataset; the outputs of these two tasks differ. Restricting supervision to a coarser prediction target can discard finer sound-type information. For example, a SPRSound Wheeze
annotation identifies a sound type also needed for ICBHI classification, although SPRSound's native binary task requires only a normal/abnormal
prediction\cite{zhang2022sprsound}. We therefore seek a joint model that learns shared acoustic attributes from available fine annotations while preserving each dataset's native prediction target.

We address this gap with LSAA (Learning Shared Acoustic Attributes), which jointly trains a shared encoder and abnormality, crackle, and wheeze predictors on ICBHI cycles and SPRSound events. Each annotation supervises only the
attributes it supports, allowing fine sound-type information to guide learning even when the native output is binary. The same model produces both native predictions through task readouts that explicitly map shared attributes to
task labels. We also use HF Lung's timed labels for auxiliary supervision and attribute evaluation, and reserve KAUH's recording-level labels for external evaluation without target-data adaptation\cite{hsu2021hflung,fraiwan2021kauh}.

Our contributions are twofold:
\begin{enumerate}
\setlength{\itemsep}{1pt}
\setlength{\parsep}{0pt}
\item We propose LSAA, a joint learning framework that connects heterogeneous
respiratory sound annotations through shared acoustic attributes, supporting
cycle-level four-class and event-level binary prediction within one model.
We further evaluate attribute transfer on HF Lung and KAUH.
\item We develop task-specific readouts over shared abnormality, crackle,
and wheeze predictions, enabling fine-grained supervision to support distinct
native outputs. Comparisons with single-source models and alternative
readouts characterize the benefits and costs of this design.
\end{enumerate}
